\sscp{Historical vowel-glide sequences *-ay, *-aw, *-uy, and *ə}{13-01-03}
\trf{13-03} shows the development of historical vowel-glide sequences
for each subgroup, as well as the development of schwa *ə.
\tsc{Taliabu} values (\crf{ch:Tal}) are included in \trf{13-03},
since these correspondences are a major part of what sets that
subgroup apart from \tsc{Sula-Buru},
\tsc{\ili{Ambon}-Seram}, and STB, and part of what connects it with \tsc{Eastern Celebic}.

\begin{table}
	\caption[Final vowel-glide sequences and *ə in each subgroup]
	{Final vowel-glide sequences and *ə in each subgroup\su{†}}\label{tab:13-03}
	\begin{threeparttable}\st{3.5pt}
	\begin{tabular}{llllll}\lsptoprule
PMP				&	*-ay										&	*-aw									&	*-uy													&	*ə\sub{1}					&	*ə\sub{2}						\\	\midrule
\tsc{BL}	&	*ay \it{(i, e, a, aʤ)}	&	*aw \it{(o, a, av)}		&	*uy \it{(i, u)}								&	*ə								&	*ə \it{(e, a, i, ə)}\\
\tsc{TB}	&	*ay \it{(a, i, e)}			&	*aw \it{(o, a)}				&	*i														&	*ə \it{(e, o)}		&	*ə \it{(e, a)}			\\
\tsc{CT}	&	*e											&	*o										&	*i														&	*ə \it{(e, o)}		&	*a									\\
\tsc{Aru}	&	*ay \it{(ay, a, i, e)}	&	*aw \it{(aw, u, o)}		&	*uy \it{(u, i,} {\0}\it{)}		&	*e, *u, *a				&	*a, *u, *i					\\
STB				&	*a											&	*a \it{(a, ə,} {\0}\it{)}	&	*uy \it{(i, u,} {\0}\it{)}&	*ə \it{(e, o)}		&	*a									\\
\ili{Banda}			&	\hr\it{a, e}						&	\hr\it{a}							&	\hr\it{u, o}									&	\hr\it{e}					&	\hr\it{a, o}				\\
\tsc{AS}	&	*a											&	*aw \it{(a, aw)}			&	*uy \it{(u, o, i}, \tsc{tr})	&	*ə \it{(e, i, o)}	&	*ə \it{(e, i, o)}		\\
\tsc{SB}	&	*a											&	*a										&	*uy \it{(u, o, i,} {\0}\it{)}	&	*e								&	\hr\it{e, i}				\\
\tsc{Tal}	&	*e											&	*o										&	*u														&	*o								&	*o									\\
		\lspbottomrule
	\end{tabular}
		\begin{tablenotes}\tnsize
			\item[†] {*ə\sub{1} is penultimate syllable; *ə\sub{2} is ultimate syllable;
								\tsc{tr} = traces requiring the full vowel-glide sequence in some etyma.}
		\end{tablenotes}
	\end{threeparttable}
\end{table}

Data for PMP *-iw are sparse; in many languages only reflexes
of *kahiw ‘tree, wood’ can be found,
but those often reflect intermediate forms (e.g. **kayu, **kau, **kai).
Or the reflexes of ubiquitous *kahiw ‘tree,
wood’ conflict with the scarce reflexes of *laRiw ‘run’ or *baliw
‘counterpart’, so there is not enough data to establish a pattern or make a
claim about regular sound correspondences.
An occurrence of one is not a pattern.

\ili{As} mentioned briefly in \srf{01-07},
we find several patterns in individual languages in the region
with regard to historical vowel-glide sequences:

\begin{enumerate}
	\item	the historical vowel-glide sequences *-ay,
				*-aw, *-uy, *-iw preserve the vowels: \it{a,
				a, u, i} and lose the glides (the coda = Blust’s “glide truncation”);
	\item	the historical vowel-glide sequences *-ay,
				*-aw, *-uy, *-iw preserve the glides as vowels: \it{i,
				u, i, u} and lose the historical vowels;
	\item	the historical vowel-glide sequences *-ay,
				*-aw, coalesce (merge, combine) to \it{e}
				and \it{o} respectively; \item	the historical vowel-glide sequences
				*-ay, *-aw, *-uy, are preserved in full,
				or have traces that require the reconstruction of the full sequence,
				or the full sequence is reflected in some words,
				but not in others (lexically specific patterns);
	\item	the historical vowel-glide sequences may be lost altogether;
	\item	different individual languages
				and lower order subgroups have a variety of different combinations
				of the above patterns,
				even within the same higher level subgroup;
	\item	several of the above patterns may be lexically specific within a language.
\end{enumerate}

So consequently the full historical vowel-glide sequences *-ay,
*-aw, *-uy, *-iw must be reconstructed to account for the variations
in the data across the region.
\ili{As} can be seen from \trf{13-03},
no splits or mergers or other innovations have been identified
that are shared by all subgroups.

\citet[264]{Blust1993} sees glide truncation (the loss of the
coda in historical vowel-glide sequences; pattern {\#}1 above)
as “one of the most distinctive characteristics of the phonological
history of CMP languages”.
In response to evidence from \citet{Nothofer1992} that glide
truncation is also found in MP languages west of the Blust line,
\citet[266]{Blust1993} clarifies that by “glide truncation” he
means the glide is lost
and the vowel retained in all four historical vowel-glide sequences,
not just in one of them.

One does not need to look very far to find significant problems with that claim.
For example, looking at patterns of sound change in \ili{Bima},
\ili{Manggarai}, \ili{Kambera} and \ili{Havu}, data provided in \citet[92]{Blust2008}
show that none of them fit the pattern of the glide being lost
and the vowel retained in all the vowel-glide sequences.
\ili{Buru} mostly follows the truncation pattern: *-ay > \it{a},
*-aw > \it{a}, *-uy > \it{u, o}.
But the reflexes of *-iw are not clear in \ili{Buru}.
For example, PMP *kahiw > \ili{Buru} \it{kau} ‘tree,
wood’, but PMP *baliw ‘dual division,
moiety [\ldots] opposite side or part; partner,
friend, enemy’
> \ili{Buru} \it{in-fali-n} ‘mother’s younger sister’.
At the higher level node of the \tsc{Sula-Buru} subgroup to which
\ili{Buru} belongs, we find: *-ay > \it{a}; *-aw > \it{a}; *-uy > \it{u,
o, i,} {\0}; *-iw > \it{i, u}.
“Truncation of the coda” does not happen fully in all vowel-glide
sequences in these \tsc{Sula-Buru} languages.
And there are similar complexities in other subgroups in the
region as shown in \trf{13-03}.
Similar variations in glide truncation are also found widely
west of the Blust line,
as in \ili{Malay} *pajay > \it{padi} ‘rice plant’,
*babuy > \it{babi} ‘pig’,
but *kahiw > \it{kayu} ‘wood,
tree’, and *kasaw > \it{kasau} ‘rafter’.
Although there are scattered individual languages that have “glide
truncation”, there are no subgroups in the region that clearly represent the
pattern of Blust’s “truncation of the coda”.

Since it is unpatterned
and not exclusive, it is difficult to see where this has subgrouping value.
\citet[72]{Blust2009} later conceded that “glide truncation,
was a product of diffusion,
since some languages that show no evidence of it are interspersed
with those that do.”

In wrapping up this section,
we note that secondary vowel sequences *a(C)i,
*a(C)u, created through the loss of historical consonants often
mirror the behaviour of *-ay
and *-aw, or *ai and *au respectively.
So *-ay/*ai/*a(C)i > \it{e},
\mbox{*-aw}/ *au/*a(C)u > \it{o}, or *au/*a(C)u > \it{ai},
or *au/*a(C)u > \it{oi}.
Yet in any given language it is fairly rare for all combinations
to show the same pattern.
It is not uncommon to have *a(C)i > \it{e} or \it{ai} in different
words in the same language.

Significantly, in languages where *-ay/*ai > \it{e} or *-aw/*au
> \it{o}, sometimes *aqi > \it{ai},
and *aqu > \it{au},
indicating the loss of *q post-dates coalescence.
For example, in \ili{Tii} (a \tsc{Timor-Babar} language of Rote),
*-aw/*-au/*-ahu > \it{o}, as shown in \qf{13-04},
which contrasts with *taqun ‘year’ > \it{teu-k} ‘year’,
showing a trace of PMP *q with raising of *aq > \it{e}.

\noindent{\wg\st{6pt}\tblr{>{\hspace{-\tabcolsep}}l<{\hspace{-\tabcolsep}}
>{\hspace{-\tabcolsep}\enspace}
llll}
\lsptoprule
%\tbx{13-04}	&	*gloss	&	PMP	&	\ili{Tii}	&	local gl.	\\	\midrule
%\endfirsthead
\tbx{13-04}	&	*gloss	&	PMP	&	\ili{Tii}	&	local gl.	\\	\midrule
%\endhead
	&	calm (water)	&	*ma-lin\tbr{aw}	&	lin\tbr{o}-k	&	calm	\\
	&	sister (m.s.)	&	*bət\tbr{aw}	&	fet\tbr{o}	&	unmarried girl, daughter	\\
	&	steal	&	*nak\tbr{aw}	&	naʔ\tbr{o}	&	steal	\\
	&	rat	&	*lab\tbr{aw}	&	laf\tbr{o}	&	mouse, rat	\\
	&	far	&	*z\tbr{au}q	&	ɗ\tbr{oo}	&	far, long (time)	\\
	&	leaf	&	*d\tbr{ahu}n	&	ɗ\tbr{oo}n	&	leaf	\\
	&	person	&	*t\tbr{au}	&	t\tbr{ou}	&	male	\\
\lspbottomrule
\end{tabular}\mbox{}\\}

The only collection of subgroups that could be linked together
on the basis of the reflexes of vowel-glide sequences in final
position are \tsc{Sula-Buru}, \tsc{\ili{Ambon}-Seram}, \ili{Banda}, and \tsc{Seram-Tanimbar-Bomberai} which
all share *-ay > \it{a}.
Yet the same change in \ili{Banda} must have come about by parallel
development and/or diffusion (see \srf{12-05}) considerably weakens
any evidence that final *-ay > \it{a} would have for such a subgroup.\footnote{
		There is some evidence that the extinct \ili{Batumerah} language,
		formerly spoken in the present day location of \ili{Ambon} city,
		may have retained *-ay as \it{ai} \citep[138]{Stresemann1927}.
		However, \citet[86--87]{Collins1983} points out several words
		in which final *a > \it{ai}
		and analyses the final \it{i} as a noun marker,
		corresponding to \it{-e} in other \tsc{\ili{Ambon}-Seram} languages.}

At best these patterns are only relevant for lower level nodes
within subgroups, and do not link subgroups together at a higher level.

